\chapter{Application Binary Interface (ABI)}
\mpitermtitleindex{ABI}
\label{sec:abi}
\label{chap:abi}

\section{Introduction}
\mpitermtitleindex{ABI!interoperability}
\label{sec:binary-interop}

The other chapters of the \MPI/ standard specify an Application Programming Interface (API) 
that defines (amongst other things) a set of opaque handle types and named constants without 
specifying their memory layout or values, respectively. 
This allows implementations to choose these according to different types of requirement. 
However, this flexibility means that different implementations are incompatible from 
the perspective of compiled applications, because the Application Binary Interface (ABI) 
is not specified.

This chapter defines an Application Binary Interface (ABI) for \MPI/,
meaning that it specifies the memory layouts of all opaque handle types, 
the values of integer constants, and other aspects of \MPI/ needed by use cases 
that require a defined ABI. 
This standard ABI for \MPI/ exists in parallel with existing implementation ABIs, 
in order to preserve backwards-compatibility of existing \MPI/ implementations.

A standard ABI for \MPI/ serves many purposes, including support for
applications compiled with one implementation of \MPI/ to be executed
with another implementation.
It is also a necessary requirement for third-party languages that intend
to interface with \MPI/ through binary symbol names, rather than direct
function calls to the C API. 
This chapter specifies the ABI of the C API of \MPI/
as well as the implications of the ABI for Fortran implementations.

\section{Implementation Requirements}

Although the ABI is designed to be portable,
there are platform and implementation designs where it cannot
be supported, or is not useful to support.
Furthermore, backwards compatibility of existing implementation
ABIs is important to some users, and \MPI/ implementations
may continue to support these.

The standard ABI is intended to support systems with the
following properties:
\begin{itemize}
    \item Dynamic shared libraries, which are loaded
          by the operating system or by the application,
          are supported.
    \item The calling convention of C functions and the
          sizes and alignment requirements of C standard
          types are known constant properties of the system;
          if the system possesses any means for changing
          these, each choice constitutes a different,
          incompatible system from the perspective of the \MPI/ ABI.
    \item Addresses can be represented as 32- or 64-bit signed
          integers and have a direct relationship with C pointers;
          segmented addressing is not supported.
\end{itemize}

%\begin{rationale}
%\end{rationale}

The following query functions are provided to allow \MPI/ applications, tools, 
and language bindings to determine whether an implementation provides ABI support
and, if so, which version of the ABI is supported.
\mpifunc{MPI\_ABI\_GET\_VERSION}
and \mpifunc{MPI\_ABI\_GET\_INFO},
can be called at any time in an \MPI/ program.
These functions must always be thread-safe, as defined in
\sectionref{sec:ei-threads}\mpitermindex{threads!thread-safe}.

\begin{mpi-binding}
    function_name("MPI_Abi_get_version")
    parameter("abi_major", "VERSION", direction="out", desc=r"ABI major version")
    parameter("abi_minor", "VERSION", direction="out", desc=r"ABI minor version")
\end{mpi-binding}

\mpifunc{MPI\_ABI\_GET\_VERSION} produces the standard ABI version, if supported.
Otherwise, the values of the major and minor version are set to $-1$.
The ABI version is independent of the \MPI/ specification version.
The major and minor version of the ABI associated with \MPIVDOTO/
are 1 and 0.

The ABI version macros \mpiconstmain{MPI\_ABI\_VERSION} and
\mpiconstmain{MPI\_ABI\_SUBVERSION} are present in the \MPI/ header and modules
so that applications can check for consistency between the compilation
environment and the properties of the implementation at runtime.

%%% Make sure that you also change these values in
%%% chap-appLang/appLang-Const.tex (C preprocessor Constants and Fortran
%%% Parameters)
\mpiconstmainindex{MPI\_ABI\_VERSION}%
\mpiconstmainindex{MPI\_ABI\_SUBVERSION}%
%%HEADER
%%SKIP
%%ENDHEADER
\begin{lstlisting}[language={[MPI]C}]
#define MPI_ABI_VERSION    1
#define MPI_ABI_SUBVERSION 0
\end{lstlisting}
%%HEADER
%%SKIP
%%ENDHEADER
\begin{lstlisting}[language={[MPI]Fortran}]
INTEGER :: MPI_ABI_VERSION, MPI_ABI_SUBVERSION
PARAMETER (MPI_ABI_VERSION    = 1)
PARAMETER (MPI_ABI_SUBVERSION = 0)
\end{lstlisting}
\noindent

Backwards-compatible changes, such as the addition of new handle types,
will increment the minor version.
Backwards-incompatible changes will increment the major version.
The addition of new functions to the \MPI/ API
does not change the ABI version.
The existing function \mpifunc{MPI\_GET\_VERSION} can be used
to query the version of the API supported and whether certain
functions are present in the \MPI/ library.

\begin{mpi-binding}
    function_name("MPI_Abi_get_info")
    parameter(
        "info",
        "INFO",
        direction="out",
        desc=r"ABI details info object (implementation-defined)",
    )
\end{mpi-binding}

Implementations may provide additional information related to the ABI.
The function \mpifunc{MPI\_ABI\_GET\_INFO} allows the user to query
this information via an info object.

The following keys are predefined for this object:
\begin{description}
\item[\infokeymain{mpi\_aint\_size}:] The size in bytes of \mpictype{MPI\_Aint}.
\item[\infokeymain{mpi\_count\_size}:] The size in bytes of \mpictype{MPI\_Count}.
\item[\infokeymain{mpi\_offset\_size}:] The size in bytes of \mpictype{MPI\_Offset}.
\end{description}

\subsection{The \texorpdfstring{\MPI/}{MPI} ABI Header File and Shared Library}
\label{subsec:c-abi-filenames}

The ABI must be implemented using a header named \code{mpi.h}.
The \MPI/ library that implements the standard ABI must be
named \code{mpi\_abi}.
The filename for this library may have a platform-specific
prefix and/or a platform-specific suffix. For Linux, for example,
\code{lib} and \code{.so} would be the default prefix and suffix.
Implementors are expected to follow platform-specific conventions
for dynamic shared library naming and versioning.
ABI-compliant implementations must not require more than \code{mpi\_abi}
or its versioned variant as the sole direct dependency of the application binary.

\begin{implementors}
If an implementation implements its own ABI definition,
it must clearly document how users employ one or the other,
such as the paths of the aforementioned files and any
other options required for their correct use.
\end{implementors}

Applications must not mix different ABIs. 
If implementations provide both the standard ABI and 
an implementation-specific ABI, applications must compile 
and link against only one of these.

The API defined in \code{mpi.h} associated with the standard ABI
does not include features of \MPI/ deprecated in \MPIIIIDOTI/ or earlier.
A full list of deprecated features can be found in Table~\ref{table:deprecated}.

\begin{rationale}
If deprecated features are included in the standard ABI, deleting
them will cause a backwards-incompatibility issue in the ABI.
Removing them from the ABI now makes it straightforward for them
to be deleted from \MPI/ in the future.
\end{rationale}

\section{The C Application Binary Interface}
\label{sec:c-abi}

\subsection{The Status Object}

The \MPI/ status object is a struct containing 8 integers: 
the three public member fields described in \sectionref{subsec:pt2pt-status}
and 5 private member fields that are reserved for implementations and 
must never be directly accessed by applications.

The \MPI/ status object is defined in C as follows:
%%HEADER
%%SKIP
%%ENDHEADER
\begin{lstlisting}[language={[MPI]C}]
typedef struct {
    int MPI_SOURCE;
    int MPI_TAG;
    int MPI_ERROR;
    int MPI_internal[5];
} MPI_Status;
\end{lstlisting}
The \MPI/ status object must use exactly eight C \ctype{int} worth of storage.

\begin{implementors}
The alignment of the status object may be less than the alignment
of a pointer or \mpictype{MPI\_Count}.  Therefore, implementatons
that store such a value in the status object must take care to access
it using a method that does not depend on alignment greater than \ctype{int}.
Such methods include \code{memcpy} and type-punning.
\end{implementors}

\begin{rationale}
This definition provides sufficient space to accommodate implementation-specific 
information and leads to memory alignment that allows efficient access to arrays
of status objects.
\end{rationale}

\subsection{Opaque Handles}
\label{subsec:abi-opaque-handles}

Handles for \MPI/ objects are defined to be incomplete struct pointers, 
which allows for C compilers to do type-checking, while also satisfying 
the existing requirements, such as equality comparison.

\begin{rationale}
Integer handles do not provide type safety, while \code{struct} or \code{union} handles 
fail to satisfy the existing API requirements.
\end{rationale}

The following handle type definitions are part of the \MPI/ ABI:

%%HEADER
%%SKIP
%%ENDHEADER
\begin{lstlisting}[language={[MPI]C}]
typedef struct MPI_ABI_Comm* MPI_Comm;
typedef struct MPI_ABI_Datatype* MPI_Datatype;
typedef struct MPI_ABI_Errhandler* MPI_Errhandler;
typedef struct MPI_ABI_File* MPI_File;
typedef struct MPI_ABI_Group* MPI_Group;
typedef struct MPI_ABI_Info* MPI_Info;
typedef struct MPI_ABI_Message* MPI_Message;
typedef struct MPI_ABI_Op* MPI_Op;
typedef struct MPI_ABI_Request* MPI_Request;
typedef struct MPI_ABI_Session* MPI_Session;
typedef struct MPI_ABI_T_cvar_handle* MPI_T_cvar_handle;
typedef struct MPI_ABI_T_enum* MPI_T_enum;
typedef struct MPI_ABI_T_event_instance* MPI_T_event_instance;
typedef struct MPI_ABI_T_event_registration* MPI_T_event_registration;
typedef struct MPI_ABI_T_pvar_handle* MPI_T_pvar_handle;
typedef struct MPI_ABI_T_pvar_session* MPI_T_pvar_session;
typedef struct MPI_ABI_Win* MPI_Win;
\end{lstlisting}

\subsection{Handle Constants}

Every handle type has at least one, and often many, 
predefined constants of that type, e.g.,
\mpiconst{MPI\_REQUEST\_NULL} for \mpictype{MPI\_Request} and 
\mpiconst{MPI\_COMM\_WORLD} for \mpictype{MPI\_Comm}.
The \MPI/ ABI defines handle constants to be compile-time constants,
which are specified as integer expressions cast to the appropriate handle type.

\begin{rationale}
Link-time constants, while convenient, are not strictly portable.
\end{rationale}

All predefined handle constants correspond to integer representations 
that are unlikely to be valid addresses, and must not be dereferenced. 
Implementations must ensure that the handles they create for the user 
are never in the range reserved for predefined handle constants. 
The \MPI/ ABI reserves values corresponding to
the integers 1 to 4095 for predefined handle constants.
The definition of these constants is described in \sectionref{sec:handle-abi}.
Handle arguments with an integer representation of zero are never valid handles.

\begin{rationale}
Many operating systems support a ``zero page'' that corresponds 
to the above address range, in which case, implementations 
will not need to do any runtime checking to ensure the above requirement is satisfied.
Ensuring that an integer representation of zero is never a legal handle
argument allows the detection of uninitialized data, which may lead to
undefined behavior.
\end{rationale}

All of the constants are specified in \sectionref{subsec:annexa-const}.

\subsection{Integer Constants}

Integer constants fall into groups, where all constants in each group must have unique values. 
In cases where integer constants are intended to be combined using bitwise logical expressions, 
there are additional requirements, specified elsewhere
(e.g. \sectionref{sec:io-filecntl-open}).
A different constraint exists for constants like \mpiconst{MPI\_ANY\_SOURCE}, 
which must be negative, because any non-negative integer may be a valid rank.

The \MPI/ ABI reserves all unused values up to 16384 for integer constants,
to allow for adding new constants without creating a noncontiguous range.
For example, implementations may define their own extensions for
\mpifunc{MPI\_COMM\_SPLIT\_TYPE} that use non-standard values of the
\mpiarg{split\_type} argument -- these integer constants must exceed 16384.

\subsection{Integer Types}

\MPI/ defines four special types of integers in C:
\mpictype{MPI\_Aint}, \mpictype{MPI\_Offset}, \mpictype{MPI\_Count},
and \mpictype{MPI\_Fint}.
The properties of \mpictype{MPI\_Aint} correspond to the C standard integer type \code{intptr\_t}.
Thus, \mpictype{MPI\_Aint} should be defined as a C \code{typedef} to \code{intptr\_t}.
In a compilation environment where \code{intptr\_t} is not available,
a type with the same properties must be used.

Essentially all filesystems relevant to \MPI/ use 64-bit addressing
so the standard ABI defines \mpictype{MPI\_Offset}
to be the C standard integer type \code{int64\_t} (or an equivalent type).
On systems where \code{intptr\_t} is 32 or 64 bits,
\mpictype{MPI\_Count} is the C standard integer type \code{int64\_t} (or an equivalent type).
% said elsewhere
%\mpictype{MPI\_Count} is required to be able to hold values of
%\mpictype{MPI\_Aint} and \mpictype{MPI\_Offset}.

\mpictype{MPI\_Fint} is discussed in \sectionref{sec:fortran-abi}.

\begin{rationale}
Fixing the size of \mpictype{MPI\_Offset} ensures the standard ABI
depends only on the address size and thus is unambiguous
on each platform.
The need for \MPI/ to support offsets greater than 64 bits
implies the possibility of a single \MPI/ file of more
than 8 exabytes in size, which is not currently practical. % revisit this assumption in 2040+
The use of 64-bit \MPI/ offsets does not prevent \MPI/
from supporting 128-bit filesystems.
\end{rationale}

\subsection{Calling Conventions and Binary Representations}

ABI compatibility 
means that the binary object code of the \MPI/ implementation, 
libraries that use \MPI/, and the main \MPI/ application 
program can be linked and executed correctly.
The type layouts, symbol names, and calling conventions of 
\MPI/ routines behave as if they have been compiled with the 
system C compiler toolchain 
(as determined, in particular, by the system C runtime library).

\begin{users}
Libraries and applications that use \MPI/ may be built with any toolchain they wish, 
as long as they adhere to these conventions when calling \MPI/ routines. 
Compiler options that change the size or layout of types, or calling conventions, 
should be avoided.
\end{users}

\section{The Fortran Application Binary Interface}
\label{sec:fortran-abi}

Fortran support for the ABI is more complicated than C,
because the sizes of \ftype{INTEGER}, \ftype{REAL},
and \ftype{DOUBLE PRECISION} can be changed with compiler options;
there is no ABI constancy even with a single Fortran compiler and platform.
This flexibility creates a difficult situation for the \MPI/ C ABI
because the functions defined in \sectionref{sec:misc-handleconvert}
and \sectionref{sec:conversion:status} in the \MPI/ C API
depend on the size of \ftype{INTEGER} via \mpictype{MPI\_Fint}.
As a result, the functions defined in \sectionref{sec:misc-handleconvert}
and \sectionref{sec:conversion:status} 
as well as \mpictype{MPI\_F08\_Status} are not part of this ABI.
Instead, to support Fortran and other use cases where obtaining an
integer associated with an \MPI/ handle is necessary, new functions are
added that do not depend on \mpictype{MPI\_Fint}.

\MPI/ applications can discover the size of Fortran types such as
\type{MPI\_INTEGER} and \type{MPI\_REAL} using \mpifunc{MPI\_TYPE\_SIZE}.
Lack of support in the implementation for optional predefined datatypes
is indicated when the type size returned is
\mpiconst{MPI\_UNDEFINED}. % the other option is 0

\begin{rationale}
Prior to the ABI, optional predefined datatypes were not present in the \MPI/ header and modules.
When absent, usage would generate compilation errors.
The standard ABI must define a value for all predefined datatypes, including the optional ones.
Therefore, the absence of optional datatypes must be detected at runtime.
\end{rationale}

\subsection{Fortran Type Registration}
\label{subsec:fortran-type-registration}

In order to decouple \MPI/ Fortran support from the rest of the implementation,
there must be a way to inform the implementation of the properties of
\MPI/ Fortran datatypes.
With this, it is possible to implement \MPI/ Fortran support on top
of the \MPI/ C implementation, without the latter having to know the
properties of the Fortran environment.

\begin{mpi-binding}
    function_name("MPI_Abi_set_fortran_info")
    parameter(
        "info",
        "INFO",
        desc=r"Fortran ABI details info object",
    )
\end{mpi-binding}

\begin{mpi-binding}
    function_name("MPI_Abi_get_fortran_info")
    parameter(
        "info",
        "INFO",
        desc=r"Fortran ABI details info object",
        direction="out"
    )
\end{mpi-binding}

\mpifunc{MPI\_ABI\_SET\_FORTRAN\_INFO} allows the application
to inform the implementation of the sizes of Fortran types
and whether or not optional types are supported by the
Fortran compiler.
Before setting this information, the application should
get this info object using \mpifunc{MPI\_ABI\_GET\_FORTRAN\_INFO}.
When \mpiconst{MPI\_INFO\_NULL} is returned, the implementation
does not know the properties of the Fortran compiler and they
must be set by the application.
%Otherwise, the application cannot set these properties.
%It should, however, check to see that the properties contained
%in the info object returned are compatible with the Fortran
%compiler and not proceed if they are incompatible.
Only the first call to \mpifunc{MPI\_ABI\_SET\_FORTRAN\_INFO}
affects the state of the \MPI/ library; all subsequent calls
will return the error code \errormain{MPI\_ERR\_ABI}.
If a call to \mpifunc{MPI\_ABI\_SET\_FORTRAN\_INFO} is not
successful, the user should call \mpifunc{MPI\_ABI\_GET\_FORTRAN\_INFO}
to determine what Fortran compiler properties were set.

The following keys are predefined for this object:
\begin{description}
\item[\infokeymain{mpi\_logical\_size}:] The size in bytes of the Fortran default \ftype{LOGICAL} kind.
\item[\infokeymain{mpi\_integer\_size}:] The size in bytes of the Fortran default \ftype{INTEGER} kind.
\item[\infokeymain{mpi\_real\_size}:] The size in bytes of the Fortran default \ftype{REAL} kind.
\item[\infokeymain{mpi\_double\_precision\_size}:] The size in bytes of the Fortran \ftype{DOUBLE PRECISION} kind.
%\item[\infokeymain{mpi\_async\_protects\_nonblocking}:] Equivalent to \mpiconst{MPI\_ASYNC\_PROTECTS\_NONBLOCKING}.
%\item[\infokeymain{mpi\_subarrays\_supported}:] Equivalent to \mpiconst{MPI\_SUBARRAYS\_SUPPORTED}.
\item[\infokeymain{mpi\_logical1\_supported}:]  (boolean) \mpidtype{MPI\_LOGICAL1} is supported.
\item[\infokeymain{mpi\_logical2\_supported}:]  (boolean) \mpidtype{MPI\_LOGICAL2} is supported.
\item[\infokeymain{mpi\_logical4\_supported}:]  (boolean) \mpidtype{MPI\_LOGICAL4} is supported.
\item[\infokeymain{mpi\_logical8\_supported}:]  (boolean) \mpidtype{MPI\_LOGICAL8} is supported.
\item[\infokeymain{mpi\_logical16\_supported}:] (boolean) \mpidtype{MPI\_LOGICAL16} is supported.
\item[\infokeymain{mpi\_integer1\_supported}:]  (boolean) \mpidtype{MPI\_INTEGER1} is supported.
\item[\infokeymain{mpi\_integer2\_supported}:]  (boolean) \mpidtype{MPI\_INTEGER2} is supported.
\item[\infokeymain{mpi\_integer4\_supported}:]  (boolean) \mpidtype{MPI\_INTEGER4} is supported.
\item[\infokeymain{mpi\_integer8\_supported}:]  (boolean) \mpidtype{MPI\_INTEGER8} is supported.
\item[\infokeymain{mpi\_integer16\_supported}:] (boolean) \mpidtype{MPI\_INTEGER16} is supported.
\item[\infokeymain{mpi\_real2\_supported}:]     (boolean) \mpidtype{MPI\_REAL2} is supported.
\item[\infokeymain{mpi\_real4\_supported}:]     (boolean) \mpidtype{MPI\_REAL4} is supported.
\item[\infokeymain{mpi\_real8\_supported}:]     (boolean) \mpidtype{MPI\_REAL8} is supported.
\item[\infokeymain{mpi\_real16\_supported}:]    (boolean) \mpidtype{MPI\_REAL16} is supported.
\item[\infokeymain{mpi\_complex4\_supported}:]  (boolean) \mpidtype{MPI\_COMPLEX4} is supported.
\item[\infokeymain{mpi\_complex8\_supported}:]  (boolean) \mpidtype{MPI\_COMPLEX8} is supported.
\item[\infokeymain{mpi\_complex16\_supported}:] (boolean) \mpidtype{MPI\_COMPLEX16} is supported.
\item[\infokeymain{mpi\_complex32\_supported}:] (boolean) \mpidtype{MPI\_COMPLEX32} is supported.
\item[\infokeymain{mpi\_double\_complex\_supported}:] (boolean) \mpidtype{MPI\_DOUBLE\_COMPLEX} is supported.
\end{description}

\begin{mpi-binding}
    function_name("MPI_Abi_set_fortran_booleans")
    parameter("logical_size", "TYPECLASS_SIZE",
              desc=r"the size of Fortran \code{LOGICAL} in bytes",
              direction="in")
    parameter("logical_true", "LOGICAL_VOID",
              desc=r"the Fortran literal value \code{.TRUE.}",
              direction="in", pointer=True)
    parameter("logical_false", "LOGICAL_VOID",
              desc=r"the Fortran literal value \code{.FALSE.}",
              direction="in", pointer=True)
\end{mpi-binding}

\begin{mpi-binding}
    function_name("MPI_Abi_get_fortran_booleans")
    parameter("logical_size", "TYPECLASS_SIZE",
              desc=r"the size of Fortran \code{LOGICAL} in bytes",
              direction="in")
    parameter("logical_true", "LOGICAL_VOID",
              desc=r"the Fortran literal value \code{.TRUE.}",
              direction="out")
    parameter("logical_false", "LOGICAL_VOID",
              desc=r"the Fortran literal value \code{.FALSE.}",
              direction="out")
    parameter("is_set", "LOGICAL",
              desc=r"flag to indicate whether the logical boolean values were set previously",
              direction="out")
\end{mpi-binding}

\mpifunc{MPI\_ABI\_SET\_FORTRAN\_BOOLEANS} allows the application
to inform the implementation of the literal values of the Fortran booleans.
Boolean literals must be passed directly so that they can be observed
by the implementation, since it may not be possible to obtain their
literal values directly.
This function has a size argument to allow it to be called before
\mpifunc{MPI\_ABI\_SET\_FORTRAN\_INFO}.
Before setting this information, the application should
check if the logical values are already set
using \mpifunc{MPI\_ABI\_GET\_FORTRAN\_BOOLEANS}.
When \mpiarg{is\_set}\mpicode{ = false}, the implementation
does not know the properties of the Fortran compiler and they
must be set by the application.
When \mpiarg{is\_set}\mpicode{ = true}, the implementation
already knows the literal values of the Fortran booleans
and they cannot be set.
As with \mpifunc{MPI\_ABI\_SET\_FORTRAN\_INFO}, only the first call
to this function affects the state of the \MPI/ library; subsequent
calls will return the error code \error{MPI\_ERR\_ABI}.

\begin{rationale}
\MPI/ does not assume that Fortran boolean literals follow the C
convention (zero is false and non-zero is true).
\end{rationale}

\subsection{The \texorpdfstring{\MPI/}{MPI} ABI Fortran Modules and Shared Library}
\label{subsec:fortran-abi-names}

The ABI must be implemented using modules named \code{mpi} and \code{mpi\_f08}.
The \MPI/ library that implements the standard ABI must be
named \code{mpifort\_abi} and follow all of the requirements stated in
\sectionref{subsec:c-abi-filenames} for naming, versioning, and dependencies.

\subsection{The Status Object}
\label{subsec:fortran-abi-status}

The \MPI/ status object is defined in Fortran as follows:
%%HEADER
%%SKIP
%%ENDHEADER
\begin{lstlisting}[language={[MPI]Fortran}]
integer, parameter :: MPI_STATUS_SIZE = 8
type, bind(C) :: Status
   integer  :: MPI_SOURCE
   integer  :: MPI_TAG
   integer  :: MPI_ERROR
   integer  :: MPI_INTERNAL(5)
end type Status
\end{lstlisting}
The \MPI/ status object must use exactly eight Fortran \ftype{INTEGER} worth of storage.

The following constants can be specified:
%%HEADER
%%SKIP
%%ENDHEADER
\begin{lstlisting}[language={[MPI]C}]
#define MPI_F_STATUS_SIZE 8
#define MPI_F_SOURCE      0
#define MPI_F_TAG         1
#define MPI_F_ERROR       2
\end{lstlisting}

\subsection{Integer Constants}

As specified elsewhere (Section~\ref{sec:misc-constants}),
constants have the same value in all languages, unless specified otherwise.

\begin{sloppy}
The constants
\mpiconst{MPI\_F\_STATUS\_IGNORE}, \mpiconst{MPI\_F\_STATUSES\_IGNORE},
\mpiconst{MPI\_F08\_STATUS\_IGNORE} and \mpiconst{MPI\_F08\_STATUSES\_IGNORE}
are not specified in the C header file.

\end{sloppy}

\begin{rationale}
Users can obtain the addresses of Fortran sentinels by passing them as
arguments to a Fortran function call that is implemented in C.
\end{rationale}

\subsection{Handle Serialization}
\mpitermtitleindex{handle serialization}
\label{sec:abi-handleserialization}

\sectionref{sec:misc-handleconvert} defines
methods for converting handles to integers of the type
\mpictype{MPI\_Fint}.
In order to provide this functionality without depending
on the behavior of the Fortran compiler,
new functions that convert handles to and from
\ctype{int} are necessary.
Because these functions depend only on C language features,
they are not referred to as C-Fortran conversion functions
but handle serialization functions.

In the C ABI, handles are pointers and therefore applications
can trivially serialize handles into the type \ctype{intptr\_t}
using a cast, but this does not support use cases where a
language integer type is narrower than this.

\begin{rationale}
While it is possible to implement this functionality outside of \MPI/,
e.g. using a lookup table or hash function,
it is possible to implement more efficiently within an \MPI/
implementation, particularly if the implementation is already
using a limited range of values for C handles.
An implementation may also store the integer associated with
a C handle, in which case the lookup is trivial.
\end{rationale}

\begin{mpi-binding}
    no_lis_binding()
    no_f08_binding()
    no_f90_binding()

    with_uppercase_index()

    function_name('MPI_Comm_toint')
    return_type("KEY_INDEX")

    parameter("comm", "COMMUNICATOR", direction="in", pointer=False)
\end{mpi-binding}

The function \cfunc{MPI\_Comm\_toint}
translates a C communicator handle into a C integer.
For all predefined handles, the integer value must be
the same as the values listed in Section~\ref{sec:lang}.
For user-defined handles, the implementation must
return the same integer for every call to this function with the same handle,
which does not conflict with the reserved range for predefined handles.
It is erroneous to call this function with an invalid handle argument.

\begin{mpi-binding}
    no_lis_binding()
    no_f08_binding()
    no_f90_binding()

    with_uppercase_index()

    function_name('MPI_Comm_fromint')
    return_type("COMMUNICATOR")

    parameter("comm", "KEY_INDEX", direction="in", pointer=False)
\end{mpi-binding}

The function \cfunc{MPI\_Comm\_fromint} translates a C integer
to the appropriate C communicator handle.
Only an integer obtained from a previous call to \cfunc{MPI\_Comm\_toint}
may be passed to this function.
It is erroneous to pass to this function an integer associated
with a handle that has been freed, disconnected, or aborted (or that was derived from a session that has been finalized).

Similar functions are provided for the other types of opaque objects.

\bindingislisttrue
\begin{mpi-binding}
    no_lis_binding()
    no_f08_binding()
    no_f90_binding()

    with_uppercase_index()

    function_name('MPI_Type_fromint')
    return_type("DATATYPE")

    parameter("datatype", "KEY_INDEX", direction="in", pointer=False)
\end{mpi-binding}

\begin{mpi-binding}
    no_lis_binding()
    no_f08_binding()
    no_f90_binding()

    with_uppercase_index()

    function_name('MPI_Type_toint')
    return_type("KEY_INDEX")

    parameter("datatype", "DATATYPE", direction="in", pointer=False)
\end{mpi-binding}

\cdeclindex{MPI\_Group}%

\begin{mpi-binding}
    no_lis_binding()
    no_f08_binding()
    no_f90_binding()

    with_uppercase_index()

    function_name('MPI_Group_fromint')
    return_type("GROUP")

    parameter("group", "KEY_INDEX", direction="in", pointer=False)
\end{mpi-binding}

\cdeclindex{MPI\_Group}%
\begin{mpi-binding}
    no_lis_binding()
    no_f08_binding()
    no_f90_binding()

    with_uppercase_index()

    function_name('MPI_Group_toint')
    return_type("KEY_INDEX")

    parameter("group", "GROUP", direction="in", pointer=False)
\end{mpi-binding}

\cdeclindex{MPI\_Request}%
\begin{mpi-binding}
    no_lis_binding()
    no_f08_binding()
    no_f90_binding()

    with_uppercase_index()

    function_name('MPI_Request_fromint')
    return_type("REQUEST")

    parameter("request", "KEY_INDEX", direction="in", pointer=False)
\end{mpi-binding}

\cdeclindex{MPI\_Request}%
\begin{mpi-binding}
    no_lis_binding()
    no_f08_binding()
    no_f90_binding()

    with_uppercase_index()

    function_name('MPI_Request_toint')
    return_type("KEY_INDEX")

    parameter("request", "REQUEST", direction="in", pointer=False)
\end{mpi-binding}

\cdeclindex{MPI\_File}%
\begin{mpi-binding}
    no_lis_binding()
    no_f08_binding()
    no_f90_binding()

    with_uppercase_index()

    function_name('MPI_File_fromint')
    return_type("FILE")

    parameter("file", "KEY_INDEX", direction="in", pointer=False)
\end{mpi-binding}

\cdeclindex{MPI\_File}%
\begin{mpi-binding}
    no_lis_binding()
    no_f08_binding()
    no_f90_binding()

    with_uppercase_index()

    function_name('MPI_File_toint')
    return_type("KEY_INDEX")

    parameter("file", "FILE", direction="in", pointer=False)
\end{mpi-binding}

\cdeclindex{MPI\_Win}%
\begin{mpi-binding}
    no_lis_binding()
    no_f08_binding()
    no_f90_binding()

    with_uppercase_index()

    function_name('MPI_Win_fromint')
    return_type("WINDOW")

    parameter("win", "KEY_INDEX", direction="in", pointer=False)
\end{mpi-binding}

\cdeclindex{MPI\_Win}%
\begin{mpi-binding}
    no_lis_binding()
    no_f08_binding()
    no_f90_binding()

    with_uppercase_index()

    function_name('MPI_Win_toint')
    return_type("KEY_INDEX")

    parameter("win", "WINDOW", direction="in", pointer=False)
\end{mpi-binding}

\cdeclindex{MPI\_Op}%
\begin{mpi-binding}
    no_lis_binding()
    no_f08_binding()
    no_f90_binding()

    with_uppercase_index()

    function_name('MPI_Op_fromint')
    return_type("OPERATION")

    parameter("op", "KEY_INDEX", direction="in", pointer=False)
\end{mpi-binding}

\cdeclindex{MPI\_Op}%
\begin{mpi-binding}
    no_lis_binding()
    no_f08_binding()
    no_f90_binding()

    with_uppercase_index()

    function_name('MPI_Op_toint')
    return_type("KEY_INDEX")

    parameter("op", "OPERATION", direction="in", pointer=False)
\end{mpi-binding}

\cdeclindex{MPI\_Info}%
\begin{mpi-binding}
    no_lis_binding()
    no_f08_binding()
    no_f90_binding()

    with_uppercase_index()

    function_name('MPI_Info_fromint')
    return_type("INFO")

    parameter("info", "KEY_INDEX", direction="in", pointer=False)
\end{mpi-binding}

\cdeclindex{MPI\_Info}%
\begin{mpi-binding}
    no_lis_binding()
    no_f08_binding()
    no_f90_binding()

    with_uppercase_index()

    function_name('MPI_Info_toint')
    return_type("KEY_INDEX")

    parameter("info", "INFO", direction="in", pointer=False)
\end{mpi-binding}

\cdeclindex{MPI\_Errhandler}%
\begin{mpi-binding}
    no_lis_binding()
    no_f08_binding()
    no_f90_binding()

    with_uppercase_index()

    function_name('MPI_Errhandler_fromint')
    return_type("ERRHANDLER")

    parameter("errhandler", "KEY_INDEX", direction="in", pointer=False)
\end{mpi-binding}

\cdeclindex{MPI\_Errhandler}%
\begin{mpi-binding}
    no_lis_binding()
    no_f08_binding()
    no_f90_binding()

    with_uppercase_index()

    function_name('MPI_Errhandler_toint')
    return_type("KEY_INDEX")

    parameter("errhandler", "ERRHANDLER", direction="in", pointer=False)
\end{mpi-binding}

\cdeclindex{MPI\_Message}%
\begin{mpi-binding}
    no_lis_binding()
    no_f08_binding()
    no_f90_binding()

    with_uppercase_index()

    function_name('MPI_Message_fromint')
    return_type("MESSAGE")

    parameter("message", "KEY_INDEX", direction="in", pointer=False)
\end{mpi-binding}

\cdeclindex{MPI\_Message}%
\begin{mpi-binding}
    no_lis_binding()
    no_f08_binding()
    no_f90_binding()

    with_uppercase_index()

    function_name('MPI_Message_toint')
    return_type("KEY_INDEX")

    parameter("message", "MESSAGE", direction="in", pointer=False)
\end{mpi-binding}

\cdeclindex{MPI\_Session}%
\begin{mpi-binding}
    no_lis_binding()
    no_f08_binding()
    no_f90_binding()

    with_uppercase_index()

    function_name('MPI_Session_fromint')
    return_type("SESSION")

    parameter("session", "KEY_INDEX", direction="in", pointer=False)
\end{mpi-binding}

\bindingislistfalse
\cdeclindex{MPI\_Session}%
\begin{mpi-binding}
    no_lis_binding()
    no_f08_binding()
    no_f90_binding()

    with_uppercase_index()

    function_name('MPI_Session_toint')
    return_type("KEY_INDEX")

    parameter("session", "SESSION", direction="in", pointer=False)
\end{mpi-binding}

Within the context of the ABI, where the layout of the status object
is known and representable directly in terms of C \ctype{int},
no serialization functionality is necessary.

\section{Handle Constants}
\label{sec:handle-abi}

Predefined handle constants are represented as integers in the range 1 to 4095.
To make it easy for implementations and tools to decode handle constants,
their values are derived from a Huffman code.
The Huffman code currently uses the lower 10 bits of the 12 bits allocated above.
In the following, we will describe constants in their binary representation,
$0b***\_***\_***\_***$, where underscores are added for readability.
We count bits from the right starting from zero; the rightmost three bits
will be denoted $2:0$.

Datatypes are the most common type of predefined handle and use the range $0b00\_10\_********$.
Bits $7:3$ identify the category and the bits $2:0$ identify the specific instances thereof.
The first set of predefined datatypes are not fixed-sized.
While \ctype{long} has a fixed size for a given platform ABI, it is not fixed across all platforms.
In contrast, types like \ctype{int32\_t} and \ftype{REAL*8} (or its Fortran standard equivalents)
have the same size on all platforms.
Fixed-sized types allow the implementation to determine the size of the datatype
from the predefined handle value itself, without a lookup table.

\begin{table}[tb]
\caption{Predefined \MPI/ datatype categories and instance values, for types without a specified size.
         All unassigned values are reserved.}
\label{table:abi:datatypes:lang_indep}
\begin{center}
\begin{small}
\begin{tabular}{|l|l|l|l|l|}
\hline
\textbf{Bits $11:8$} & \textbf{Bits $7:3$} & \textbf{Bits $2:0$} & \textbf{Category} & \textbf{Instance} \\
\hline
\textbf{0010} & \textbf{00000} & & Language-independent types & \\
              &                & \textbf{000} & & \mpidtype{MPI\_DATATYPE\_NULL} \\
              &                & \textbf{001} & & \mpidtype{MPI\_AINT} \\
              &                & \textbf{010} & & \mpidtype{MPI\_COUNT} \\
              &                & \textbf{011} & & \mpidtype{MPI\_OFFSET} \\
              &                & \textbf{111} & & \mpidtype{MPI\_PACKED} \\
\hline
\end{tabular}
\end{small}
\end{center}
\end{table}

\begin{table}[tb]
\caption{Predefined \MPI/ datatype categories and instance values, for types without a specified size.
         All unassigned values are reserved.}
\label{table:abi:datatypes:c_integers}
\begin{center}
\begin{tabular}{|l|l|l|l|l|}
\hline
\textbf{Bits $11:8$} & \textbf{Bits $7:3$} & \textbf{Bits $2:0$} & \textbf{Category} & \textbf{Instance} \\
\hline
\textbf{0010} & \textbf{00001} & & C integer types & \\
              &                & \textbf{000} & & \mpidtype{MPI\_SHORT} \\
              &                & \textbf{001} & & \mpidtype{MPI\_INT} \\
              &                & \textbf{010} & & \mpidtype{MPI\_LONG} \\
              &                & \textbf{011} & & \mpidtype{MPI\_LONG\_LONG} \\
              &                & \textbf{100} & & \mpidtype{MPI\_UNSIGNED\_SHORT} \\
              &                & \textbf{101} & & \mpidtype{MPI\_UNSIGNED} \\
              &                & \textbf{110} & & \mpidtype{MPI\_UNSIGNED\_LONG} \\
              &                & \textbf{111} & & \mpidtype{MPI\_UNSIGNED\_LONG\_LONG} \\
\hline
\end{tabular}
\end{center}
\end{table}

\begin{table}[ht!]
\caption{Predefined \MPI/ datatype categories and instance values, for types without a specified width.
         All unassigned values are reserved.}
\label{table:abi:datatypes:c_floats}
\begin{center}
\begin{small}
\begin{tabular}{|l|l|l|l|l|}
\hline
\textbf{Bits $11:8$} & \textbf{Bits $7:3$} & \textbf{Bits $2:0$} & \textbf{Category} & \textbf{Instance} \\
\hline
\textbf{0010} & \textbf{00010} & & C/C++ floating-           & \\
              &                & &               point types & \\
              &                & \textbf{000} & & \mpidtype{MPI\_FLOAT} \\
              &                & \textbf{001} & & \mpidtype{MPI\_C\_FLOAT\_COMPLEX} \\
              &                & \textbf{010} & & \mpidtype{MPI\_CXX\_FLOAT\_COMPLEX} \\
              &                & \textbf{100} & & \mpidtype{MPI\_DOUBLE} \\
              &                & \textbf{101} & & \mpidtype{MPI\_C\_DOUBLE\_COMPLEX} \\
              &                & \textbf{110} & & \mpidtype{MPI\_CXX\_DOUBLE\_COMPLEX} \\
\hline
\end{tabular}
\end{small}
\end{center}
\end{table}

\begin{table}[ht!]
\caption{Predefined \MPI/ datatype categories and instance values, for types without a specified width.
         All unassigned values are reserved.}
\label{table:abi:datatypes:fortran}
\begin{center}
\begin{tabular}{|l|l|l|l|l|}
\hline
\textbf{Bits $11:8$} & \textbf{Bits $7:3$} & \textbf{Bits $2:0$} & \textbf{Category} & \textbf{Instance} \\
\hline
\textbf{0010} & \textbf{00011} & & Fortran types & \\
              &                & \textbf{000} & & \mpidtype{MPI\_LOGICAL} \\
              &                & \textbf{001} & & \mpidtype{MPI\_INTEGER} \\
              &                & \textbf{010} & & \mpidtype{MPI\_REAL} \\
              &                & \textbf{011} & & \mpidtype{MPI\_COMPLEX} \\
              &                & \textbf{100} & & \mpidtype{MPI\_DOUBLE\_PRECISION} \\
              &                & \textbf{101} & & \mpidtype{MPI\_DOUBLE\_COMPLEX} \\
              &                & \textbf{110} & & \mpidtype{MPI\_CHARACTER} \\
\hline
\end{tabular}
\end{center}
\end{table}

\begin{table}[tb]
\caption{Predefined \MPI/ datatype categories and instance values, for types without a specified width.
         All unassigned values are reserved.}
\label{table:abi:datatypes:long_double}
\begin{center}
\begin{small}
\begin{tabular}{|l|l|l|l|l|}
\hline
\textbf{Bits $11:8$} & \textbf{Bits $7:3$} & \textbf{Bits $2:0$} & \textbf{Category} & \textbf{Instance} \\
\hline
\textbf{0010} & \textbf{00100} & & Long double & \\
              &                & &       types & \\
              &                & \textbf{000} & & \mpidtype{MPI\_LONG\_DOUBLE} \\
              &                & \textbf{100} & & \mpidtype{MPI\_C\_LONG\_DOUBLE\_COMPLEX} \\
              &                & \textbf{101} & & \mpidtype{MPI\_CXX\_LONG\_DOUBLE\_COMPLEX} \\
\hline
\end{tabular}
\end{small}
\end{center}
\end{table}

\begin{table}[tb]
\caption{Predefined \MPI/ datatype categories and instance values, for types without a specified width.
         All unassigned values are reserved.}
\label{table:abi:datatypes:pair_types}
\begin{center}
\begin{tabular}{|l|l|l|l|l|}
\hline
\textbf{Bits $11:8$} & \textbf{Bits $7:3$} & \textbf{Bits $2:0$} & \textbf{Category} & \textbf{Instance} \\
\hline
\textbf{0010} & \textbf{00101} & & C pair types & \\
              &                & \textbf{000} & & \mpidtype{MPI\_FLOAT\_INT} \\
              &                & \textbf{001} & & \mpidtype{MPI\_DOUBLE\_INT} \\
              &                & \textbf{010} & & \mpidtype{MPI\_LONG\_INT} \\
              &                & \textbf{011} & & \mpidtype{MPI\_2INT} \\
              &                & \textbf{100} & & \mpidtype{MPI\_SHORT\_INT} \\
              &                & \textbf{101} & & \mpidtype{MPI\_LONG\_DOUBLE\_INT} \\
\hline
\textbf{0010} & \textbf{00110} & & Fortran pair types & \\
              &                & \textbf{000} & & \mpidtype{MPI\_2REAL} \\
              &                & \textbf{001} & & \mpidtype{MPI\_2DOUBLE\_PRECISION} \\
              &                & \textbf{010} & & \mpidtype{MPI\_2INTEGER} \\
\hline
\end{tabular}
\end{center}
\end{table}

\begin{table}[tb]
\caption{Predefined \MPI/ datatype categories and instance values, for types without a specified width.
         All unassigned values are reserved.}
\label{table:abi:datatypes:other}
\begin{center}
\begin{tabular}{|l|l|l|l|l|}
\hline
\textbf{Bits $11:8$} & \textbf{Bits $7:3$} & \textbf{Bits $2:0$} & \textbf{Category} & \textbf{Instance} \\
\hline
\textbf{0010} & \textbf{00111} & & Other C/C++ types & \\
              &                & \textbf{000} & & \mpidtype{MPI\_C\_BOOL} \\
              &                & \textbf{001} & & \mpidtype{MPI\_CXX\_BOOL} \\
              &                & \textbf{100} & & \mpidtype{MPI\_WCHAR} \\
\hline
\end{tabular}
\end{center}
\end{table}

The Huffman code for fixed-size types encodes the base-2 logarithm of
the type size in bits $5:3$.  Implemenations can identify datatypes
with fixed-size using bit $6$.

% scheme
% C/C++ (0b01_***)
% 0b000: "MPI_INT(n)_T"
% 0b001: "MPI_UINT(n)_T"
% 0b010: "<float (n)b>"
% 0b011: (size=1) ? "MPI_CHAR" : "<C complex 2x(n/2)b>"
% 0b100: (size=1) ? "MPI_SIGNED_CHAR" : "reserved datatype"
% 0b101: (size=1) ? "MPI_UNSIGNED_CHAR" : "reserved datatype"
% 0b110: (size=2) ? "<C++23 bfloat16_t>" : "reserved datatype"
% 0b111: (size=1) ? "MPI_BYTE" : "<C++ complex 2x(n/2)b>"
% Fortran (0b11_***)
% 0b000: "MPI_INTEGER(n)"
% 0b001: "MPI_LOGICAL(n)"
% 0b010: "MPI_REAL(n)"

\begin{table}[tb]
\caption{Predefined \MPI/ datatype categories and instance values, for types with a specified width.
         All unassigned values are reserved.}
\label{table:abi:datatypes:c_fixed}
\begin{center}
\begin{tabular}{|l|l|l|l|l|}
\hline
\textbf{Bits $11:8$} & \textbf{Bits $7:3$} & \textbf{Bits $2:0$} & \textbf{Category} & \textbf{Instance} \\
\hline
\textbf{0010} & \textbf{01000} & & 1-byte C/C++ types & \\
              &                & \textbf{000} & & \mpidtype{MPI\_INT8\_T} \\
              &                & \textbf{001} & & \mpidtype{MPI\_UINT8\_T} \\
              &                & \textbf{011} & & \mpidtype{MPI\_CHAR} \\
              &                & \textbf{100} & & \mpidtype{MPI\_SIGNED\_CHAR} \\
              &                & \textbf{101} & & \mpidtype{MPI\_UNSIGNED\_CHAR} \\
              &                & \textbf{111} & & \mpidtype{MPI\_BYTE} \\
\hline
\textbf{0010} & \textbf{01001} & & 2-byte C/C++ types & \\
              &                & \textbf{000} & & \mpidtype{MPI\_INT16\_T} \\
              &                & \textbf{001} & & \mpidtype{MPI\_UINT16\_T} \\
\hline
\textbf{0010} & \textbf{01010} & & 4-byte C/C++ types & \\
              &                & \textbf{000} & & \mpidtype{MPI\_INT32\_T} \\
              &                & \textbf{001} & & \mpidtype{MPI\_UINT32\_T} \\
\hline
\textbf{0010} & \textbf{01011} & & 8-byte C/C++ types & \\
              &                & \textbf{000} & & \mpidtype{MPI\_INT64\_T} \\
              &                & \textbf{001} & & \mpidtype{MPI\_UINT64\_T} \\

\hline
\end{tabular}
\end{center}
\end{table}

\begin{table}[tb]
\caption{Predefined \MPI/ datatype categories and instance values, for types with a specified width.
         All unassigned values are reserved.}
\label{table:abi:datatypes:fortran_fixed}
\begin{center}
\begin{tabular}{|l|l|l|l|l|}
\hline
\textbf{Bits $11:8$} & \textbf{Bits $7:3$} & \textbf{Bits $2:0$} & \textbf{Category} & \textbf{Instance} \\
\hline
\textbf{0010} & \textbf{11000} & & 1-byte Fortran types & \\
              &                & \textbf{000} & & \mpidtype{MPI\_LOGICAL1} \\
              &                & \textbf{001} & & \mpidtype{MPI\_INTEGER1} \\
\hline
\textbf{0010} & \textbf{11001} & & 2-byte Fortran types & \\
              &                & \textbf{000} & & \mpidtype{MPI\_LOGICAL2} \\
              &                & \textbf{001} & & \mpidtype{MPI\_INTEGER2} \\
              &                & \textbf{010} & & \mpidtype{MPI\_REAL2} \\
\hline
\textbf{0010} & \textbf{11010} & & 4-byte Fortran types & \\
              &                & \textbf{000} & & \mpidtype{MPI\_LOGICAL4} \\
              &                & \textbf{001} & & \mpidtype{MPI\_INTEGER4} \\
              &                & \textbf{010} & & \mpidtype{MPI\_REAL4} \\
              &                & \textbf{011} & & \mpidtype{MPI\_COMPLEX4} \\
\hline
\textbf{0010} & \textbf{11011} & & 8-byte Fortran types & \\
              &                & \textbf{000} & & \mpidtype{MPI\_LOGICAL8} \\
              &                & \textbf{001} & & \mpidtype{MPI\_INTEGER8} \\
              &                & \textbf{010} & & \mpidtype{MPI\_REAL8} \\
              &                & \textbf{011} & & \mpidtype{MPI\_COMPLEX8} \\
\hline
\textbf{0010} & \textbf{11100} & & 16-byte Fortran types & \\
              &                & \textbf{000} & & \mpidtype{MPI\_LOGICAL16} \\
              &                & \textbf{001} & & \mpidtype{MPI\_INTEGER16} \\
              &                & \textbf{010} & & \mpidtype{MPI\_REAL16} \\
              &                & \textbf{011} & & \mpidtype{MPI\_COMPLEX16} \\
\hline
\textbf{0010} & \textbf{11101} & & 32-byte Fortran types & \\
              &                & \textbf{011} & & \mpidtype{MPI\_COMPLEX32} \\

\hline
\end{tabular}
\end{center}
\end{table}

Reduction operators use the range $0b00\_00\_001\_*****$ and the predefined values
fall into four categories: arithmetic, bit-wise, logical, and other.

All other predefined handles use the range $0b00\_01\_********$.
In general, predefined null handles have all zero bits other than those required to detect
the handle type.

\begin{table}[tb]
\caption{Predefined \mpictype{MPI\_Op} categories and instance values.
         All unassigned values are reserved.}
\label{table:abi:ops}
\begin{center}
\begin{tabular}{|l|l|l|l|l|}
\hline
\textbf{Bits $11:5$} & \textbf{Bits $4:3$} & \textbf{Bits $2:0$} & \textbf{Category} & \textbf{Instance} \\
\hline
\textbf{0000001} & \textbf{00} & & Arithmetic operations & \\
                 &             & \textbf{000} & & \mpiconst{MPI\_OP\_NULL} \\
                 &             & \textbf{001} & & \mpiconst{MPI\_SUM} \\
                 &             & \textbf{010} & & \mpiconst{MPI\_MIN} \\
                 &             & \textbf{011} & & \mpiconst{MPI\_MAX} \\
                 &             & \textbf{100} & & \mpiconst{MPI\_PROD} \\
\hline
\textbf{0000001} & \textbf{01} & & Bit operations & \\
                 &             & \textbf{000} & & \mpiconst{MPI\_BAND} \\
                 &             & \textbf{001} & & \mpiconst{MPI\_BOR} \\
                 &             & \textbf{010} & & \mpiconst{MPI\_BXOR} \\
\hline
\textbf{0000001} & \textbf{10} & & Logical operations & \\
                 &             & \textbf{000} & & \mpiconst{MPI\_LAND} \\
                 &             & \textbf{001} & & \mpiconst{MPI\_LOR} \\
                 &             & \textbf{010} & & \mpiconst{MPI\_LXOR} \\
\hline
\textbf{0000001} & \textbf{11} & & Other operations & \\
                 &             & \textbf{000} & & \mpiconst{MPI\_MINLOC} \\
                 &             & \textbf{001} & & \mpiconst{MPI\_MAXLOC} \\
                 &             & \textbf{100} & & \mpiconst{MPI\_REPLACE} \\
                 &             & \textbf{101} & & \mpiconst{MPI\_NO\_OP} \\

\hline
\end{tabular}
\end{center}
\end{table}

\begin{table}[tb]
\caption{Predefined \MPI/ handle categories and instance values.
         All unassigned values are reserved.}
\label{table:abi:handles}
\begin{center}
\begin{tabular}{|l|l|l|l|l|}
\hline
\textbf{Bits $11:8$} & \textbf{Bits $7:3$} & \textbf{Bits $2:0$} & \textbf{Category} & \textbf{Instance} \\
\hline
\textbf{0001} & \textbf{00000} & & Communicators & \\
              &                & \textbf{000} & & \mpiconst{MPI\_COMM\_NULL} \\
              &                & \textbf{001} & & \mpiconst{MPI\_COMM\_WORLD} \\
              &                & \textbf{010} & & \mpiconst{MPI\_COMM\_SELF} \\

\textbf{0001} & \textbf{00001} & & Group & \\
              &                & \textbf{000} & & \mpiconst{MPI\_GROUP\_NULL} \\
              &                & \textbf{001} & & \mpiconst{MPI\_GROUP\_EMPTY} \\

\textbf{0001} & \textbf{00010} & & Windows & \\
              &                & \textbf{000} & & \mpiconst{MPI\_WIN\_NULL} \\

\textbf{0001} & \textbf{00011} & & Files & \\
              &                & \textbf{000} & & \mpiconst{MPI\_FILE\_NULL} \\

\textbf{0001} & \textbf{00100} & & Sessions & \\
              &                & \textbf{000} & & \mpiconst{MPI\_SESSION\_NULL} \\

\textbf{0001} & \textbf{00101} & & Messages & \\
              &                & \textbf{000} & & \mpiconst{MPI\_MESSAGE\_NULL} \\
              &                & \textbf{001} & & \mpiconst{MPI\_MESSAGE\_NO\_PROC} \\

\textbf{0001} & \textbf{00110} & & Info & \\
              &                & \textbf{000} & & \mpiconst{MPI\_INFO\_NULL} \\
              &                & \textbf{001} & & \mpiconst{MPI\_INFO\_ENV} \\

\textbf{0001} & \textbf{01000} & & Errhandlers & \\
              &                & \textbf{000} & & \mpiconst{MPI\_ERRHANDLER\_NULL} \\
              &                & \textbf{001} & & \mpiconst{MPI\_ERRORS\_ARE\_FATAL} \\
              &                & \textbf{010} & & \mpiconst{MPI\_ERRORS\_RETURN} \\
              &                & \textbf{011} & & \mpiconst{MPI\_ERRORS\_ABORT} \\

\textbf{0001} & \textbf{10000} & & Requests & \\
              &                & \textbf{000} & & \mpiconst{MPI\_REQUEST\_NULL} \\

\hline
\end{tabular}
\end{center}
\end{table}

